% ============================================================
%  LAB / LEARNING JOURNAL
%  Build with latexmk (LaTeX Workshop's default recipe).
%  Save the file and a PDF should appear; or press Ctrl+Alt+B.
% ============================================================
\documentclass[11pt]{article}

% --- Page geometry: comfortable margins for reading/note-taking ---
\usepackage[a4paper, margin=1in]{geometry}   % use letterpaper if you're in the US

% --- Fonts / encoding (modern, sensible defaults) ---
\usepackage[T1]{fontenc}
\usepackage{lmodern}

% --- Math: amsmath is the workhorse; amssymb/amsthm add symbols & theorem tools ---
\usepackage{amsmath, amssymb, amsthm}

% --- Units: siunitx handles numbers + units correctly, e.g. \SI{9.81}{m/s^2} ---
\usepackage{siunitx}

% --- Figures (the actual \includegraphics example below is commented out) ---
\usepackage{graphicx}

% --- Code listings: no external tools needed (unlike minted, which wants shell-escape) ---
\usepackage{listings}
\usepackage{xcolor}
\usepackage{float}
\setcounter{tocdepth}{3}
\lstset{
  basicstyle=\ttfamily\small,
  keywordstyle=\color{blue!70!black},
  commentstyle=\color{green!50!black},
  stringstyle=\color{purple!70!black},
  numbers=left, numberstyle=\tiny\color{gray},
  breaklines=true, frame=single, language=Python
}

% --- hyperref makes the table of contents and references clickable. Load it LAST. ---
\usepackage[colorlinks=true, linkcolor=blue!50!black, urlcolor=blue!50!black]{hyperref}

\usepackage{tabularx}
\usepackage{booktabs}

\usepackage{etoc}
\usepackage{csquotes}
% ------------------------------------------------------------
%  Custom tools
% ------------------------------------------------------------

% A top-level heading with no date — for standing sections like the papers list.
\newcommand{\standing}[1]{\clearpage\section{#1}}

% A paper: a subsection that starts on its own page (for the Papers section).
\newcommand{\paper}[1]{\clearpage\subsection{#1}}

% An "entry" is just a section with a date built into the heading,
% so each one shows up in the table of contents automatically.
\newcommand{\entry}[2]{\clearpage\section[#2]{#1\quad\textnormal{\textit{#2}}}}
%   usage: \entry{2026-06-25}{What I learned about limits}

% A sub-entry: a subsection that nests under the current entry.
\newcommand{\subentry}[1]{\subsection{#1}}
%   usage: \subentry{A smaller topic within today's entry}

\newcommand{\subsubentry}[1]{\subsubsection{#1}}

% A box for results worth flagging. Definitions/claims you want to find later.
\theoremstyle{definition}
\newtheorem{keyresult}{Key result}
\newtheorem{definition}{Definition}
\newtheorem*{thoughts}{Thoughts}
\newtheorem*{paraphrase}{In my words}
\newtheorem{question}{Open question}   
\newtheorem*{answer}{Answer}           
\newtheorem*{connection}{Connects to}
% ------------------------------------------------------------
\title{Gaussian Process Multi-fidelity Modeling of Near-Earth Galactic Cosmic Ray Spectra}
\author{Ian Silber}                      % put your name here if you like
\date{\today}                  % the cover date; entries carry their own dates

\begin{document}
\maketitle
\tableofcontents               % grows automatically as you add entries
\vspace{1em}\hrule\vspace{1em}

% ============================================================
%  ENTRIES  —  copy the block below to start a new one
% ============================================================
% ------------------------------------------------------------
\entry{2026-06-25}{Radial Basis Functions: fixed and parametric}
\subentry{The general idea}

\entry{2026-06-25}{Kriging: A special form of RBF}

\begin{definition}[Symmetric Positive Definite (SPD) Matrix]
  A real $n \times n$ matrix $A$ is symmetric positive definite if it is
  symmetric ($A = A^T$) and
  \[
    x^T A x > 0 \quad \forall\, x.
  \]
  $A$ has a unique Cholesky factorization $A = R^TR$, where $R$ is upper triangular with positive diagonal elements.
  \\
  $A$ has positive eigenvalues
\end{definition}

\begin{definition}[Maximum Likelihood Estimation (MLE)]
    We can either find the probability of the data given parameters, or find
    the parameters given the data. In the case of our Kriging models, we aim to estimate the parameters based on
    our limited data in order to create our surrogate.
\end{definition}

\begin{definition}[Stochastic Process]
    A model of a system that evolves over time or space in a random unpredictable manner.
\end{definition}

\begin{definition}[Covariance]
A measure of the correlation between two or more sets of random variables
\begin{align}
\operatorname{cov}(X,Y) &= E[(X - \mu_x)(Y - \mu_y)] \\
                        &= E[XY] - \mu_x\mu_y
\label{eq:covariance}
\end{align}
where $\mu_x$ and $\mu_y$ are the means of $X$ and $Y$ and $E$ is the expectation. To make this make more sense,
we can derive the more intuitive correlation from covariance by:
\begin{equation}
\operatorname{cor}(X,Y) = \frac{\operatorname{cov}(X,Y)}{\sigma_x\sigma_y}
\label{eq:correlation}
\end{equation}
(TODO) I still need to work on understanding what this means exactly.
\end{definition}

\entry{2026-06-29}{Multi-Fidelity Analysis: Using multiple datasets for better approximations}
\subentry{The general idea}
In the event that we have multiple datasets describing a process,
it is useful to combine data from an expensive function and a cheap approximation of 
said function. We want to sample the expensive function at a select few points that overlap with
the sampling of the cheap function ($X_e \subset X_c$). Then we can try to correct points from $Y_c$
with:
\begin{equation}
  y_e = Z_py_c + Z_d
  \label{eq:multi-fidelity}
\end{equation}

\subentry{Co-Kriging}
From Forrester:
\begin{displayquote}[Forrester et al., 2008]
  Here we use the auto-regressive model of Kennedy and O'Hagan (2000), which
  assumes that
  \begin{equation*}
    \operatorname{cov}\{Y_e(x_i),\, Y_c(x) \mid Y_c(x_i)\} = 0
    \quad \forall\, x \neq x_i.
  \end{equation*}
  This means that no more can be learnt about $Y_e(x_i)$ from the cheaper code if
  the value of the expensive function at $x_i$ is known (this is known as a
  \emph{Markov property} which, in essence, says we assume that the expensive
  simulation is correct and any inaccuracies lie wholly in the cheaper
  simulation).
\end{displayquote}
As in regular Kriging, we treat our expensive function as a stochastic process $Z_e$,
\begin{equation}
  Z_e(\mathbf{x}) = \rho Z_c(\mathbf{x}) + Z_d(\mathbf{x})
  \label{eq:co-kriging}
\end{equation}

Basically, we are modeling the expensive function as our cheap function multiplied by some
correction factor $\rho$ plus a separate difference model $Z_d$ that estimates the difference
between our cheap and expensive function at a point.

It helps me to see the final product and then dissect the parts. To use our co-kriging model to estimate
a point $x^{(n_e + 1)}$ we evaluate:

\begin{equation}
  \hat{y}_e\!\left(x^{(n_e+1)}\right)
    = \hat{\mu} + \mathbf{c}^T \mathbf{C}^{-1}\left(\mathbf{y} - \mathbf{1}\hat{\mu}\right)
  \label{eq:co-kriging-evaluate}
\end{equation}

Where:
\begin{description}
  \item[$\hat{\mu}$] the estimated global mean of the co-Kriging process
  \item[$\mathbf{c}$] a column vector of correlations between the new point $x^{(n_e+1)}$ and the observed points
  $$\mathbf{c} = 
  \begin{pmatrix} 
    \hat{\rho}\hat{\sigma}^2\psi_c(\mathbf{X}_c,x^{(n+1)}) \\ 
    \hat{\rho}^2\hat{\sigma}^2\psi_c(\mathbf{X_e},x^{(n+1)}) + \hat{\sigma}^2\psi_d(\mathbf{X_e},x^{(n+1)})
  \end{pmatrix}$$
  \item[$\mathbf{C}$] the correlation matrix of the observed points
  $$\mathbf{C} = 
  \begin{pmatrix}
    \sigma^2\mathbf{\Psi}_c(\mathbf{X}_c,\mathbf{X}_c) &\rho\sigma^2\mathbf{\Psi}_c(\mathbf{X}_c,\mathbf{X}_e) \\ 
    \rho\sigma_c^2\mathbf{\Psi}_c(\mathbf{X}_e,\mathbf{X}_c) & \rho^2\sigma^2\mathbf{\Psi}_c(\mathbf{X}_e,\mathbf{X}_e) + \sigma^2\mathbf{\Psi}_d(\mathbf{X_e},\mathbf{X}_e)
  \end{pmatrix}$$
  \item[$\mathbf{y}$] the vector of observed responses (cheap and expensive stacked)
  \item[$\mathbf{1}$] a column vector of ones, of the same length as $\mathbf{y}$
  \item[$\psi$ or $\mathbf{\Psi}$] is a column vector or matrix calculated at each point with. Optimal theta values are found 
  for each dimension and for both cheap points and expensive points ($\psi_d$)
    $$\psi^{(i)} = \exp{(-\sum_{j=1}^{k}\hat{\theta}_{j(\text{c or d})}\|x_j^{(n+1)} - x_j^{(i)}\|)}$$
\end{description}

\subentry{Co-Kriging -- sampling plan revisited}
From Forrester et al. (2007), our sampled points are:
\begin{equation}
  \mathbf{X} = \begin{pmatrix} \mathbf{X_c} \\ \mathbf{X_e} \end{pmatrix} = (x_c^{(1)},...,x_c^{(n_c)},x_e^{(1)},...,x_e^{(n_e)})^T
  \label{eq:co-kriging-sampling}
\end{equation}
The difficulty is that $$X_e \subset X_c$$
meaning we need to generate a random latin hypercube for our cheap points, and then pick the optimal subset of expensive points.

\begin{thoughts}
I think the easiest way to do this in code is to generate the optimal sampling plan for the subset, and then add points.
The problem is that we need to be able to iterate over our sample set pairing one expensive sample point to one cheap sample point.
I am trying to come up with an elegant solution that doesn't involve sorting the sample points in some messy way. Maybe I should just add
points based on the Morris-Mitchell criterion one at a time? How can I pick a new point that maximizes distance from existing points?
OK, I changed my mind, I think the optimal way to do this is finding a subset. If I try adding new points there is simply no way to ensure I have found the best plan.
If I take a subset, I run into a similar problem. Since my expensive points must be a subset, I am not able to perturb the points, which means they will likely not be
the best plan.

\end{thoughts}

Here's what I came up with, I believe it is an accurate implementation of the exchange algorithm set out in Forrester et al. (2007) \S 4. 

\begin{lstlisting}


def find_subset(cube: np.ndarray, n: int, q: int = 2, iter: int = 5) -> np.ndarray:

    rng = np.random.default_rng()
    X_tracker = {
        "phi": math.inf,
        "X": np.empty([])
    }

    for _ in range(iter):
        # restart to avoid local minima
        X_best   = rng.choice(cube, size=n, replace=False)
        phi_best = mm_phiq(X_best, q)

        if phi_best < X_tracker["phi"]:
            X_tracker["phi"] = phi_best
            X_tracker["X"]   = X_best

        for i in range(n):
            # iterate over points in our subset
            
            for point in cube:
                # try each point from larger sample set
                X_test    = X_tracker["X"].copy()
                X_test[i] = point
                phi_best  = mm_phiq(X_test, q)

                if phi_best < X_tracker["phi"]:
                    X_tracker["phi"] = phi_best
                    X_tracker["X"]   = X_test

        
    return X_tracker["X"]


def align_subset(X_c: np.ndarray, X_e: np.ndarray) -> np.ndarray:

    if X_c.shape[1] != X_e.shape[1]:
        raise ValueError("X_c and X_e must have same number of dimensions")
    
    rows_to_delete = []
    for row in X_e:
        rows_to_delete.append(np.where((X_c == row).all(axis=1))[0])

    rows_to_delete = np.concatenate(rows_to_delete)
    
    X_c   = np.delete(X_c, rows_to_delete, axis=0)
    X_tot = np.concatenate((X_e, X_c), axis=0)

    return X_tot
    
\end{lstlisting}

 \texttt{find\_subset} works to find the best subsample of our original sampling plan. Then, \texttt{align\_subset} deletes all the
subsamples from our original sampling plan before prepending them back. This helps us pair the expensive
sample points with their cheap counterparts later in the co-Kriging process.

\standing{Papers}
\etocsetnexttocdepth{subsection}
\localtableofcontents

\paper{Guzman et al. (2026) --- Toward Explaining GCR Modulation in a Dynamic Heliosheath}

\subsubentry{Definition Table}
\begin{description}
  \item[Heliosheath (HS)] is the region of the heliosphere between the termination shock and the heliopause.
  \item[Unipolar region (UHS)] 
  \item[Sectored region (SHS)] 
  \item[Spatial effects] relate to persistent plasma properties affecting particle transport over large areas. Examples 
  include the heliospheric current sheet.
  \item[Temporal effects] are transient modulation events and solar cycle variations advected outward
    from the Sun, generally lagging 1-2 yr from their creation to their arrival
    at the HS.
\end{description}

\subsubentry{Abstract}
GCR transport appears to be dominated by large, transient modulation structures
advected from the Sun, while the role of spatially dictated diffusion is likely secondary.
The stated goal of the paper is to find out the approximate relative contribution of each of these two 
factors. hello
\begin{question}\label{q:spatial}
What is spatially dictated diffusion?
\end{question}

\begin{answer}[to Question~\ref{q:spatial}]
Spatially dictated diffusion refers to diffusion caused by persistent features like the 
heliospheric current sheet, which is a large-scale structure in the heliosphere 
that can influence the transport of GCRs. 
\end{answer}

\subsubentry{Analytic Heliosphere Model}


\paper{Puzzoni et al. (2026) --- A new web-based tool for GCR modulation in the heliosphere}
\subsubentry{Abstract}
This GCR tool is based on the Gleeson and Axford force field approximation solution to the parker transport equation and 
incorporates solar modulation via the correlation between GCRs and the mean sunspot number. 
This is a highly simplified model chosen for its fast evaluation speed and ease of use.
This model is also validated against lots of empirical data sets including Voyager 1 and 2, and New Horizons.

\begin{definition}[Parker Transport Equation]
The Parker transport equation is a fundamental model in space physics. It describes how
energetic charged particles move through the solar system's magnetized and turbulent
plasma. 
\end{definition}
\subsubentry{Model equations}
As mentioned, this model is an implementation of the Gleeson and Axford solution to the parker transport equation.
This solution is found when cosmic ray streaming flux is equal to zero.
\begin{equation}
  \vec{S} = -\mathbf{\kappa} \cdot \nabla{f} - \frac{\vec{U}}{3}\frac{\partial f}{\partial \ln{p}} = 0
  \label{eq:parker-transport}
\end{equation}

Where:
\begin{description}
  \item[$\vec{S}$] Cosmic ray streaming flux. Set to zero for high energy GCRs
  \item[$\vec{U}$] Solar wind speed
  \item[$\frac{\partial f}{\partial \ln{p}}$] The momentum gradient. This term links the spatial and energy evolution of particles
  \item[$\nabla{f}$] Gradient of the phase space distribution function
\end{description}

Again, this solution is only valid when cosmic ray streaming flux $\vec{S} = 0$, which is only true
for high energy GCRs. Additionally, the model assumes a spherically symmetric, steady state heliosphere
centered on the Sun, which is of course, not realistic.

Solving Eq. 8, one obtains for $f$
\begin{equation}
  f(r, E) = f(R_{HP}, E + \phi)
  \label{eq:phase-space-distribution-function}
\end{equation}
Where:
\begin{description}
  \item[$R_{HP}$] is the radial distance to the Heliopause
  \item[$E$] is the kinetic energy of the GCRs
  \item[$\phi$] is the modulation potential
\end{description}

$\phi$ is found by

\begin{equation}
  \phi = \int_{r}^{R_{HP}}\frac{U(r')}{3g(r')}dr'
  \label{eq:phi}
\end{equation}

In the above, $U$ is still solar wind speed, and $g(r')$ is the radial dependence of the spatial diffusion
coefficient, which has the form $k(r,p) = g(r)wp$, $p$ again being momentum, and $w$ being particle velocity.
This modulation potential is where all the tunable parameters exist.

\paper{M. C. Kennedy and A. O'Hagan (2000) --- Predicting the output from a complex computer code}

\paper{Forrester et al. (2007) --- Multi-fidelity optimization via surrogate modelling}

\end{document}